% Abstract

%\renewcommand{\abstractname}{Abstract} % Uncomment to change the name of the abstract

\pdfbookmark[1]{Abstract}{Abstract} % Bookmark name visible in a PDF viewer

\begingroup
\let\clearpage\relax
\let\cleardoublepage\relax
\let\cleardoublepage\relax

\chapter*{Abstract}
\addcontentsline{toc}{chapter}{\textsc{Abstract}}

In various fields, a compelling challenge is to influence systems to align with desired behaviors. Formal methods offer ways to address this challenge by utilizing abstractions of real-world systems to compute the solutions. This thesis focuses on the control of biological systems, which contributes to solving cell reprogramming.

Biological systems are modeled using a variety of approaches, including Boolean networks (BNs), which, thanks to their simplicity and expressiveness, effectively capture interactions at the cellular level. While control problems for fully specified BNs have been extensively studied, the corresponding problem has received much less attention when the Boolean functions are only partially known.

Partially specified Boolean networks (\PBN{s}) accommodate model uncertainty by representing unknown Boolean update functions with function symbols. However, this introduces complexities in the control problem, including variations in the state transition graph between varying interpretations and optimization challenges in minimizing perturbations while ensuring target reachability.

This thesis deals with these challenges and brings the control problem to the partially specified setting of Boolean networks. The state perturbations are used to exercise control on the models since their implementation in practice can be realized via gene knockouts or overexpressions. The main difference in PSBN control is that, instead of computing only a minimal perturbation, the methods compute a set of possible perturbations together with their robustness with respect to model uncertainty.

We present and solve two flavors of the \PBN{} control problem -- source-target control and phenotype control. Two methods were developed to deal with the source-target problem -- a semi-symbolic parallel method, which uses just one-step perturbations and a symbolic method, which is able to employ one-step, temporary and permanent perturbations. The weakness of source-target control lies in the necessity to know the target attractor beforehand, which can be highly inconvenient for \PBN{}. This weakness is mitigated by the phenotype control, where only characteristic traits of the target are sufficient to be specified and also, the control works globally regardless of the source state. A method based on permanent perturbations with symbolic state space exploration was developed to solve this problem. Finally, software encompassing all these methods is described, and the method applicability is presented on real-world case studies.

\endgroup			

\vfill